\chapter{Introdução}
ESCREVER INTROOO

\section{Público-Alvo}
Artistas independentes que realizam trabalhos personalizados 
sob encomenda e necessitam de uma plataforma para organizar 
clientes e acompanhar pedidos.

\section{Objetivo do Sistema}
Desenvolver uma aplicação cujo objetivo é auxiliar artistas 
independentes que trabalham sob comissão a organizar e gerenciar
seus pedidos. O sistema permite o controle completo de clientes,
solicitações, prazos e valores, proporcionando maior organização
e profissionalismo.

\section{Requisitos Funcionais}
\begin{itemize}
    \item Cadastrar, editar, visualizar e excluir pedidos (CRUD);
    \item Cadastrar e gerenciar clientes;
    \item Adicionar observações aos pedidos;
    \item Registrar especificações (referências, medidas, cores, detalhes);
    \item Definir valor e prazo após análise do administrador.
\end{itemize}


\section{Disciplinas e Tecnologias Envolvidas}

\subsection{Disciplinas}

\textbf{Algoritmos e Programação I:}
A disciplina de Algoritmos e Programação I forneceu a base estrutural 
para o desenvolvimento do sistema, aplicando conceitos de lógica de 
programação, estruturas de controle e organização de código na linguagem 
Java. Esses conhecimentos foram essenciais para a implementação das 4 
operações básicas que compõem o núcleo de um sistema CRUD funcional.

\textbf{Banco de Dados:}
Os conceitos de modelagem e sistema de gerenciamento de banco de dados 
relacional (SGBDR) foram fundamentais para o desenvolvimento da aplicação. 
A disciplina auxiliou na definição das entidades, relacionamentos e 
normalização de tabelas, conceitos fundamentados no modelo 
entidade-relacionamento (ER) proposto por~\cite{chen1976}.

\textbf{Conceitos da Computação:}
A matéria de Conceitos da Computação introduziu fundamentos teóricos de 
criptografia, os quais serão aplicados no sistema para o armazenamento 
seguro de senhas dos usuários, uma vez que o sistema conta com dois perfis 
de acesso: o Administrador (ADM), responsável pelo gerenciamento das 
solicitações, e o Usuário (USR), o cliente que irá efetuar a solicitação 
do serviço.

\textbf{Usabilidade, Interface e UX:}
Os estudos de Usabilidade, Interface e UX, baseados na metodologia de 
Design Thinking proposta por~\cite{brown2010}, orientaram o processo de 
ideação, viabilidade e prototipação da interface e usabilidade do sistema. 
A abordagem centrada no usuário foi aplicada para moldar uma aplicação 
intuitiva e alinhada às necessidades do público-alvo, utilizando o Figma 
como ferramenta de prototipação visual.

\subsection{Tecnologias}

\textbf{Java:}
Java é uma linguagem de programação amplamente estabelecida no mercado, 
ocupando altas posições em rankings relevantes como o TIOBE Index~\cite{tiobe} 
e o Stack Overflow Developer Survey~\cite{stackoverflow_survey}. 
Foi escolhida como linguagem principal do projeto devido à familiaridade 
da equipe com a tecnologia e ao interesse em aprofundar os conhecimentos 
adquiridos ao longo do curso~\cite{java}.

\textbf{MySQL:}
Falta usar para poder formular~\cite{mysql}.

\textbf{Figma:}
O Figma é uma plataforma de design colaborativo baseada em nuvem~\cite{figma}.
Foi escolhido como ferramenta de prototipação por ser amplamente utilizado 
no mercado e por ser a principal ferramenta de design adotada pelo curso.

\textbf{draw.io:}
O draw.io é uma ferramenta online de diagramação, modelagem visual 
e criação de fluxogramas~\cite{drawio}. Foi selecionada devido à 
sua alta portabilidade, baixa curva de aprendizado e suporte 
facilitado à colaboração entre os membros da equipe.

\textbf{Git/Github:}
O Git é um sistema de controle de versão distribuído e o GitHub é 
uma plataforma baseada em nuvem que hospeda repositórios Git. Essas 
duas ferramentas foram utilizadas para o compartilhamento, controle 
e versionamento do projeto como um todo, cujo repositório está 
disponível em~\cite{repo}.

\textbf{LaTeX/abnTeX2:}
O LaTeX é um sistema de composição de documentos amplamente utilizado 
no meio acadêmico e científico, reconhecido pela qualidade tipográfica 
e pela capacidade de produzir documentos complexos de forma estruturada. 
Em conjunto com a classe abnTeX2~\cite{abntex2}, que adapta o LaTeX às 
normas da Associação Brasileira de Normas Técnicas (ABNT)~\cite{abnt14724}, 
foi utilizado para a produção da documentação acadêmica deste projeto, 
garantindo conformidade com as normas vigentes~\cite{abnt6023}.